% !TeX root = ../../main.tex
% ============================================================
% Feature Matrix Construction — prose + code
% CSC491 Textbook, Chapter: Machine Learning
%
% Preamble requirements:
%     \usepackage{minted}
%     \usepackage{caption}
%
% Compile with:
%     xelatex --shell-escape main.tex
% ============================================================

\begin{figure}[ht]
\centering
\begin{minted}[
  style=default, 
  linenos,
  frame=lines,
  framesep=3mm,
  fontsize=\small,
  baselinestretch=1.1,
]{python}
fig, axes = plt.subplots(1, 3, figsize=(15, 4))

for ax, (label, yt, yp) in zip(axes, [
    ("A (d=0)", y_test_a, pred_a),
]):
    cm = confusion_matrix(yt, yp)
    sns.heatmap(cm, annot=True, fmt=",d", cmap="Blues", ax=ax,
                xticklabels=["Down", "Up"], yticklabels=["Down", "Up"])
    ax.set_title(f"Pipeline {label}")
    ax.set_xlabel("Predicted")
    ax.set_ylabel("Actual")

plt.suptitle("Confusion Matrices", fontsize=14, weight="bold", y=1.02)
plt.tight_layout()
plt.show()
\end{minted}
\caption{Plotting confusion matrices for each pipeline. Each matrix shows the counts of true positives, false positives, true negatives, and false negatives for the model's predictions compared to the actual labels.
requires \texttt{seaborn} and \texttt{matplotlib} for visualization.}
\label{lst:Python-Confusion-Matrix}
\end{figure}